\chapter{エネルギー、曲率、放射}

\lectureoverview{ニュートン模型は、二つの制限を取り除くと格段に多くのことを明らかにする。第一に、共動銀河の全エネルギーは厳密にゼロである必要はない。そのエネルギーを任意に保つと、一般相対論が空間曲率として再解釈する\(a^{-2}\)の項が導入される。第二に、宇宙の源は無圧力の物質である必要はない。光子の波長はスケール因子とともに伸びるので、放射は\(a^{-4}\)として薄まり、法則\(a(t)\propto t^{1/2}\)を導く。したがって、物質と放射の両方を含む宇宙は、小さい\(a\)では放射優勢であり、より大きい\(a\)では物質優勢である。その過程で、我々は局所的なニュートン的運動と大域的な相対論的幾何学を区別し、一般相対論なしには決着できない問題を特定する。}

\section{ニュートンが見ることができるもの}

模型を拡張する前に、なぜ以前のニュートン的計算がそもそも成功したのかを問うべきである。物理的な宇宙は一般相対論によって記述され、その空間幾何学は大域的にユークリッド的である必要はない。それでも、任意の滑らかな曲がった空間は、十分に小さな近傍では近似的に平坦である。地球表面の小さな部分は平面的に見える。同じように、宇宙の曲率半径よりはるかに小さい領域は平坦として扱うことができる。

``ここでいう「小さい」''とは、曲率スケールに対して相対的である。その部分は依然として何十億光年にも及ぶかもしれない。その中では、隣接する銀河は、それらの相対的な固有速度が非相対論的であり、重力場が弱いとき、ニュートン力学で扱うことができる。したがって、ニュートン的計算はアインシュタインの理論の競合相手ではない。それは、より完全な記述の局所的で低速の極限である。

ただちに警告がある。放射はゆっくり動かない。光子、および相対論的である間のニュートリノは、通常のニュートン的粒子の気体として扱うことはできない。ニュートン的推論は膨張方程式を組織し続けるが、その右辺の源は静止質量密度からエネルギー密度へと拡大されなければならない。ここで近似は自らの境界を明らかにし始める。

\section{スケール因子の再確認}

銀河の一様な分布を取り、平均的な宇宙流れとともに動く格子に座標を付ける。固定された共動座標\(x\)にある銀河は、選ばれた原点から物理的距離
\begin{equation}
D(t)=a(t)x
\label{eq:c02-distance}
\end{equation}
を持つ。スケール因子\(a(t)\)は座標の隔たりを物理的な隔たりに変換する。

\(a\)の絶対値は慣習的である。すべての共動座標を定数倍で再スケールすると、同じ物理的距離は\(a\)の逆スケーリングによって表すことができる。したがって、\(a\)も、一つの座標セルに割り当てられた質量も、それぞれ独立に不変ではない。スケール因子の比と膨張率は不変である：
\begin{equation}
\frac{a(t_2)}{a(t_1)},
\qquad
H(t)\equiv\frac{\dot a}{a}.
\label{eq:c02-hubble}
\end{equation}
実際、固定された定数\(\alpha\)に対して\(a'=\alpha a\)ならば、\(\dot a'/a'=\dot a/a\)である。

\(\nu_m\)を、共動座標体積の単位体積に含まれる質量とする。その数値は格子に新たなラベルを付けると変化するが、通常の物質が保存されるときには時間的に一定である。共動単位立方体は物理的体積\(a^3\)を持つので、物質密度は
\begin{equation}
\boxed{\rho_m(a)=\frac{\nu_m}{a^3}.}
\label{eq:c02-matter-density}
\end{equation}
である。この\(a^{-3}\)の薄まりは単に体積の増大である：共動箱の中の静止エネルギーは固定されたままで、その体積が増大する。

実際の銀河は格子に堅く釘付けにされているわけではない。それらは群れの中で軌道を描き、銀河団へと落ち込み、平均的な流れに対して小さな速度を持つ。スケール因子は大きなスケールでの平滑化された運動を記述する。そのような局所的なずれは\emph{固有速度}と呼ばれる。

\begin{classroomqa}
\audiencequestion{すべての銀河は、他のすべての銀河から厳密に動径方向に遠ざからなければならないのか。}
\lecturerresponse{そのような形を持つのは、平均化されたハッブル流れだけである。銀河の実際の運動は、共通の膨張と、近傍の構造によって引き起こされる固有速度の和である。この区別は、渦が重なった川の流れのようなものである。渦は局所的には重要だが、大規模な流れを消し去ることはない。\lecturetimestamp{00:24:47--00:26:33}}
\end{classroomqa}

ニュートン的全エネルギーがゼロの無圧力物質に対して、前回の講義では
\begin{equation}
H^2=\frac{8\pi G}{3}\rho_m
=\frac{8\pi G\nu_m}{3a^3}.
\label{eq:c02-zero-energy-friedmann}
\end{equation}
が得られた。全体的な格子の規格化は任意であるから、一時的に\(8\pi G\nu_m/3=1\)となる単位を選んでよい。膨張する枝では、
\begin{equation}
\frac{\dot a}{a}=\frac{1}{a^{3/2}},
\qquad
\dot a=\frac{1}{\sqrt a}.
\end{equation}
である。導関数を反転して積分すると、
\begin{equation}
\frac{dt}{da}=\sqrt a,
\qquad
t-t_B=\frac{2}{3}a^{3/2},
\end{equation}
が得られる。ここで\(t_B\)は、この理想化された解が\(a=0\)に達する時刻である。したがって、
\begin{equation}
\boxed{a(t)\propto (t-t_B)^{2/3}.}
\label{eq:c02-matter-law}
\end{equation}
である。この臨界物質のみの模型では、宇宙は永遠に膨張する。その膨張は\(\dot a\propto a^{-1/2}\)なので減速するが、\(\dot a\)は\(a\)のいかなる有限の値でもゼロにはならない。

\begin{figure}[ht]
\centering
\includegraphics[width=0.68\textwidth]{lecture_02_figure_02.png}
\caption{復習板には\(\dot a=1/\sqrt a\)と枠で囲まれた物質支配の法則\(a\propto t^{2/3}\)が記録されている。中間の代数計算は一部が隠されており、式~\eqref{eq:c02-zero-energy-friedmann}--\eqref{eq:c02-matter-law}で補われる。}
\label{fig:c02-matter-recap}
\lectureframe{00:29:20.838}
\end{figure}

\begin{classroomqa}
\audiencequestion{なぜ負の平方根は捨てられたのか。}
\lecturerresponse{それは方程式によって捨てられたのではない。式~\eqref{eq:c02-zero-energy-friedmann}が定めるのは\(H^2\)であり、\(H\)の符号ではない。正の根は膨張を表し、負の根は収縮を表す。惑星の重力場の中の岩は、同じエネルギー方程式に従いながら外向きにも内向きにも運動しうる。どちらの枝を取るかは初期条件が選ぶ。我々は議論中の宇宙に対応する膨張の枝を選んだ。\lecturetimestamp{00:23:09--00:24:46}}
\end{classroomqa}

次の模型でも依然として暗黒エネルギーは省かれる。この省略は意図的である。すなわち、現在の加速期を支配する成分を加える前に、まず古い物質と放射の枠組みを理解するのである。これに対して暗黒物質はすでに物質項に含まれている。非相対論的な暗黒物質もまた\(a^{-3}\)として薄まるからである。

\begin{classroomqa}
\audiencequestion{暗黒エネルギーは今日の支配的成分ではないのか。}
\lecturerresponse{はい。ここで展開した二成分の話は、現在の加速の発見以前の標準的な教育用宇宙論である。暗黒物質は物質に属し、暗黒エネルギーは別に導入されるべき第三の成分である。古い模型から始めることで、追加される各項が何を変えるのかが見えてくる。\lecturetimestamp{00:28:02--00:28:45}}
\end{classroomqa}

\section{エネルギー定数の復元}

ゼロエネルギーの選択は特別な初期条件である。これを除くには、一様な球の境界上にある質量\(m\)の銀河に戻る。その共動半径を\(x\)とすると、
\begin{equation}
D=ax,
\qquad
\dot D=\dot a x.
\label{eq:c02-comoving-particle}
\end{equation}
ニュートンの球殻定理によれば、球の内部の質量\(M(D)\)だけが動径方向の加速度に寄与し、外側の球殻は正味の力を及ぼさない。銀河の保存されたニュートンエネルギーは
\begin{equation}
\frac{1}{2}m\dot D^{\,2}-\frac{GM(D)m}{D}=E.
\label{eq:c02-newton-energy}
\end{equation}
因子\(m\)は両方の項に属する。板上で急いで書くときには省略しやすいが、これを保つことで、最終的な定数が単位試験質量あたりのエネルギーであることが明確になる。

一様性により、代表銀河を\(x=1\)に選べる。すると\(D=a\)、\(\dot D=\dot a\)であり、球は
\begin{equation}
V_{\rm sph}=\frac{4\pi}{3}a^3,
\qquad
M(a)=\rho_m V_{\rm sph}=\frac{4\pi}{3}\rho_m a^3.
\label{eq:c02-enclosed-mass}
\end{equation}
板では\(V\)が体積にも、傍注では動径速度にも使われている。我々は曖昧さのない記号\(V_{\rm sph}\)と\(\dot D\)を用いることにする。

\begin{figure}[ht]
\centering
\includegraphics[width=0.86\textwidth]{lecture_02_figure_03.png}
\caption{囲まれた質量の球と非ゼロエネルギーのフリードマン方程式。見えるラベルには\(V_{\rm sph}=4\pi a^3/3\)、中心質量\(M\)、\(x=1\)、\(D=a\)、\(\dot D=\dot a\)が含まれる。方程式の一部は講師によって隠されており、以下で代数的に再構成される。}
\label{fig:c02-energy-board}
\lectureframe{00:37:59.876}
\end{figure}

\begin{figure}[ht]
\centering
\begin{tikzpicture}[scale=1.05,>=Latex]
  \draw[thick] (0,0) circle (2.6);
  \fill (0,0) circle (2.2pt) node[below left=2pt] {\(M(a)\)};
  \draw[thick,->] (0,0) -- (2.58,0);
  \fill (2.6,0) circle (2.2pt) node[right=4pt] {galaxy \(m\)};
  \node[above] at (1.30,0.04) {\(D=ax\)};
  \node[below] at (1.30,-0.04) {\(x=1\Rightarrow D=a\)};
  \draw[->,blue!60!black,thick] (2.63,0.18) -- (3.55,0.72)
    node[right] {\(\dot D=\dot a\)};
  \node at (0,-3.05) {\(M(a)=\frac{4\pi}{3}\rho_m a^3\)};
\end{tikzpicture}
\caption{\(x=1\)におけるニュートン球の明快な再構成。外側の球殻は打ち消し合うので、境界の銀河は囲まれた質量にのみ応答する。}
\label{fig:c02-sphere-redraw}
\end{figure}

\(D=a\)を式~\eqref{eq:c02-newton-energy}に代入し、\(m/2\)で割り、定数を
\begin{equation}
C\equiv\frac{2E}{m}.
\end{equation}
と定義する。すると
\begin{equation}
\dot a^{\,2}-\frac{2GM(a)}{a}=C.
\label{eq:c02-energy-a}
\end{equation}
が得られる。\(a^2\)で割り、式~\eqref{eq:c02-enclosed-mass}を用いると
\begin{equation}
\left(\frac{\dot a}{a}\right)^2-
\frac{8\pi G}{3}\rho_m=\frac{C}{a^2},
\label{eq:c02-energy-form}
\end{equation}
あるいは等価に
\begin{equation}
\boxed{H^2=\frac{8\pi G}{3}\rho_m+\frac{C}{a^2}.}
\label{eq:c02-friedmann-C}
\end{equation}
が得られる。これは圧力のない物質に対する第一フリードマン方程式であり、ニュートンのエネルギー規約で書かれている。新しい\(a^{-2}\)の項は、代表銀河が脱出エネルギーより上、ちょうど、あるいは下で放たれたかどうかを記憶している。

\(\rho_m=\nu_m/a^3\)を代入すると、
\begin{equation}
A_m\equiv\frac{8\pi G\nu_m}{3}>0,
\end{equation}
と定義するのが便利であり、その結果
\begin{equation}
H^2=\frac{A_m}{a^3}+\frac{C}{a^2},
\qquad
\dot a^{\,2}=\frac{A_m}{a}+C.
\label{eq:c02-energy-branches}
\end{equation}
となる。この方程式は、直ちに積分するよりも極限によって読むほうがより多くのことを語る。

\subsection{正のエネルギー}

\(C>0\)ならば、小さな\(a\)では物質項\(A_m/a^3\)が支配する。したがって初期の膨張は物質の法則\(a\propto t^{2/3}\)を保つ。しかし大きな\(a\)では、\(H^2\)において\(C/a^2\)の項が支配し、等価に
\begin{equation}
\dot a^{\,2}\longrightarrow C.
\end{equation}
膨張は一定速度に近づく。
\begin{equation}
a(t)\sim\sqrt C\,t.
\label{eq:c02-positive-late}
\end{equation}
これは、脱出速度より速くリンゴを投げることに相当する宇宙論的な場合である。重力は絶えずその速度を減らすが、遠方ではリンゴは非ゼロの漸近速度を保つ。

\subsection{ゼロエネルギー}

\(C=0\)に対して、式~\eqref{eq:c02-energy-branches}は臨界解\(a\propto t^{2/3}\)を返す。それは\(\dot a\)がゼロに近づきながら無限に膨張する。これはちょうど脱出エネルギーの場合である。

\subsection{負のエネルギー}

\(C=-|C|<0\)ならば、
\begin{equation}
H^2=\frac{A_m}{a^3}-\frac{|C|}{a^2}
=\frac{A_m-|C|a}{a^3}.
\label{eq:c02-negative-energy}
\end{equation}
スケール因子は
\begin{equation}
a_{\max}=\frac{A_m}{|C|},
\label{eq:c02-amax}
\end{equation}
を超えられない。右辺がそこでゼロに達するからである。膨張の枝は\(\dot a=0\)とともに\(a_{\max}\)に到達し、その続きは収縮の枝である。ニュートンの描像では、これは脱出速度より遅く放たれたリンゴである。すなわち、それは上昇し、折り返し、落下する。

\begin{classroomqa}
\audiencequestion{右辺が\(H^2\)に等しく、したがって負にはなり得ないのに、なぜモデルは再収縮し得るのか。}
\lecturerresponse{決して負にはならない。膨張の間、右辺は\(a_{\max}\)でゼロまで減少する。そこでは\(H=0\)。その後、解は正の平方根から負の平方根へと変わり、\(\dot a<0\)として許された領域を逆向きにたどる。誘惑されるが誤った継続は、この表式をゼロを通り過ぎさせるであろう。それは不可能である。ゼロは境界であり、実数の\(H\)が通り抜けて継続できる値ではない。\lecturetimestamp{00:48:43--00:53:03}}
\end{classroomqa}

したがって、三つのニュートン的エネルギーは三つの定性的な歴史を与える。\(C<0\)に対しては再収縮、\(C=0\)に対しては臨界的脱出、そして\(C>0\)に対しては\(\dot a\)が漸近的に一定のまま永遠に膨張する場合である。驚くべき次の段階は、一般相対論がこれらの同じ場合に幾何学的な意味を割り当てることである。

\section{エネルギーから曲率へ}

方程式~\eqref{eq:c02-friedmann-C}のエネルギー項を幾何学的な側へ移し、
\begin{equation}
k\equiv-C.
\label{eq:c02-k-definition}
\end{equation}
と定義する。すると
\begin{equation}
\boxed{H^2+\frac{k}{a^2}=\frac{8\pi G}{3}\rho.}
\label{eq:c02-friedmann-k}
\end{equation}
黒板では\(C\)、\(K\)、\(\kappa\)の間で交互に用いられ、項を等号の反対側へ移す際に符号が吸収される。方程式~\eqref{eq:c02-k-definition}が我々の規約を固定する。正のニュートン的エネルギーは\(k<0\)に対応し、ゼロエネルギーは\(k=0\)に対応し、負のニュートン的エネルギーは\(k>0\)に対応する。

\begin{figure}[ht]
\centering
\includegraphics[width=0.76\textwidth]{lecture_02_figure_04.png}
\caption{フリードマン方程式を、左辺に幾何学、右辺に源として整理したもの。この枠組みはこの概念的な分割を支えるが、完全なテンソル形式のアインシュタイン方程式を表示しているわけではない。}
\label{fig:c02-geometry-source}
\lectureframe{00:55:22.810}
\end{figure}

この整理はアインシュタイン方程式の縮小版である。左辺は幾何学を記述する。すなわち\(H\)による膨張と、\(k/a^2\)による空間曲率である。右辺は物質的な源を記述する。ニュートン力学ではその源は質量密度であった。相対論ではそれはエネルギーと運動量であり、完全には応力エネルギー・テンソルによって表される。一様等方な宇宙に対して、第一フリードマン方程式は全エネルギー密度\(\rho\)に依存する。

通常の単位を保つならば、質量密度はその静止エネルギー\(\rho_{\rm mass}c^2\)を通じて寄与する。以下では\(c=1\)と置くので、質量とエネルギーは同じ単位を用いる。そのとき記号\(\rho\)は静止質量、運動エネルギー、放射、およびその他の任意の一様なエネルギー成分を含み得る。

\begin{figure}[ht]
\centering
\includegraphics[width=0.88\textwidth]{lecture_02_figure_05.png}
\caption{フリードマン方程式の並置された黒板形式。\(\rho\)を用いる曲率形式と、\(\rho_m=\nu_m/a^3\)の後に得られる物質形式。黒板上で記法は変わるが、方程式~\eqref{eq:c02-k-definition}はここでの符号規約を固定したままである。}
\label{fig:c02-curvature-comparison}
\lectureframe{00:57:52.659}
\end{figure}

幾何学的解釈は大域的である。相対論的モデルでは、\(k=0\)は平坦な空間的断面を記述し、\(k>0\)は正に曲がった球面的断面を記述し、\(k<0\)は負に曲がった双曲的断面を記述する。ニュートン的導出は、これらの大域的幾何学を持たずに同じ発展方程式を再現した。したがってその非ゼロのエネルギー定数は、曲率の驚くほど正確な模倣ではあるが、曲がった空間そのものの導出ではない。

\begin{classroomqa}
\audiencequestion{ニュートン力学と一般相対論の幾何学が異なるのに、なぜ同じ膨張方程式を生み出し得るのか。}
\lecturerresponse{一般相対論は平坦な場合を含み、局所的にはすべての滑らかな幾何学は近似的に平坦である。さらに驚くべきことに、ニュートン的エネルギー定数は、相対論的な空間曲率が占めるまさにその場所で発展方程式に入る。動力学の式は一致するが、大域的解釈は一致しない。ユークリッド的なニュートン空間が密かに球面や双曲空間になったわけではない。三つの幾何学は一般相対論の内部で構成されなければならない。\lecturetimestamp{01:20:59--01:24:20}}
\end{classroomqa}

\section{膨張する箱の中の放射}

ここで宇宙の内容を変える。一つの共動単位立方体を考える。その物理的体積は\(a^3\)である。非相対論的物質では、その箱の中の粒子数は固定され、各粒子の静止エネルギーはほぼ固定されるので、\(\rho_m\propto a^{-3}\)が得られる。放射には別の効果がある。各光子は宇宙の膨張に伴ってエネルギーを失う。

波長\(\lambda\)の光子はエネルギー
\begin{equation}
E_\gamma=\frac{hc}{\lambda},
\qquad
E_\gamma\propto\lambda^{-1}.
\label{eq:c02-photon-energy}
\end{equation}
を持つ。一様膨張の下では、自由に伝播する波長はスケール因子とともに伸びる：
\begin{equation}
\frac{\lambda(t_2)}{\lambda(t_1)}
=\frac{a(t_2)}{a(t_1)},
\qquad
\lambda\propto a.
\label{eq:c02-redshift}
\end{equation}
これが宇宙論的赤方偏移である。これは次を意味する：
\begin{equation}
E_\gamma\propto a^{-1}.
\label{eq:c02-photon-redshift}
\end{equation}
共動箱内の光子数が保存されるなら、その全放射エネルギーは\(a^{-1}\)として減少する。これをその物理的体積\(a^3\)で割ると
\begin{equation}
\boxed{\rho_r(a)=\frac{\nu_r}{a^4}.}
\label{eq:c02-radiation-density}
\end{equation}
が得られる。ここで\(\nu_r\)は自由に流れる放射密度の一定振幅である。放射は物質よりも\(a\)の逆冪を1つ多く持つ：空間的希釈による3つの冪と、各光子の赤方偏移による1つである。

ギター弦のアナロジーは有用な直観を与える。弦に適合するモードは整数個の半波長を持つ。弦がゆっくり長くされると、その整数は連続的に変化しえない；代わりに波長が新しい長さに調整される。同様に、ゆっくり膨張する環上の波では、節の数は断熱不変量であり、したがってその波長は円周に従う。一般の波はそのようなモードに分解できる。

これは膨張時空におけるマクスウェル方程式を解くことの代用ではない。それは方程式~\eqref{eq:c02-redshift}の背後にある制御された直観である：ゆっくりした幾何学的膨張はモードラベルを保ちながら物理的波長を変える。

\begin{classroomqa}
\audiencequestion{この議論には導体または反射箱が必要だろうか？物理的な壁のない空の膨張空間では何が起こるのか？}
\lecturerresponse{壁は帳簿上の道具であり、宇宙が導体に囲まれているという主張ではない。完全な証明は膨張計量における波の伝播から従う。当面の議論では、自由に伝播する波長が\(a\)とともにスケールするという、観測されまた相対論的に導かれた事実を用いる。箱は光子と物理的体積を数えるのに役立つ；それは宇宙論的赤方偏移を引き起こすのではない。\lecturetimestamp{01:02:21--01:04:37}}
\end{classroomqa}

\begin{classroomqa}
\audiencequestion{なぜ波長は固定された物理的長さを保つのではなく、格子に追随すべきなのか？}
\lecturerresponse{円周が断熱的に変化する環上のモードを考える。その整数の節の数は連続的にずれえないので、許される波長は環とともに変化しなければならない。一般の波をフーリエモードに展開すると、モードごとに同じ結論が得られる。これは必要な断熱的直観を与える；厳密な定式化は曲がった時空における場の理論に属する。\lecturetimestamp{01:34:18--01:38:33}}
\end{classroomqa}

\subsection{放射優勢解}

今しばらく曲率を省き、方程式~\eqref{eq:c02-radiation-density}をフリードマン方程式に代入すると：
\begin{equation}
H^2=\frac{8\pi G\nu_r}{3a^4}
\equiv\frac{A_r}{a^4},
\qquad A_r>0.
\label{eq:c02-radiation-friedmann}
\end{equation}
膨張する枝では、
\begin{equation}
\frac{\dot a}{a}=\frac{\sqrt{A_r}}{a^2},
\qquad
\dot a=\frac{\sqrt{A_r}}{a}.
\end{equation}
したがって
\begin{equation}
a\,da=\sqrt{A_r}\,dt,
\qquad
\frac{a^2}{2}=\sqrt{A_r}(t-t_B),
\end{equation}
そして
\begin{equation}
\boxed{a(t)\propto(t-t_B)^{1/2}.}
\label{eq:c02-radiation-law}
\end{equation}
指数\(1/2\)と\(2/3\)の区別は見た目には控えめだが、物理的には決定的である。初期宇宙における膨張率、温度、反応の凍結、地平線の大きさはすべて、どちらの成分が優勢かによって決まる。

\section{物質と放射の宇宙}

別々に保存される物質と放射については、それらの密度は加算される：
\begin{equation}
\rho(a)=\rho_m(a)+\rho_r(a).
\end{equation}
この比較のために曲率を無視すると、
\begin{equation}
\boxed{H^2=\frac{C_m}{a^3}+\frac{C_r}{a^4},}
\label{eq:c02-mixed-friedmann}
\end{equation}
ここで\(C_m,C_r>0\)である。板書ではまず2つの係数が\(C_1\)と\(C\)と記される；2番目の記号は口頭で別個の定数に訂正される。方程式~\eqref{eq:c02-mixed-friedmann}の記法はその訂正を明示している。

\begin{figure}[ht]
\centering
\includegraphics[width=0.74\textwidth]{lecture_02_figure_06.png}
\caption{板上の物質と放射の混合方程式。書かれた2番目の係数は後に放射に固有の定数へ訂正され、方程式~\eqref{eq:c02-mixed-friedmann}では\(C_r\)と記される。}
\label{fig:c02-mixed-board}
\lectureframe{01:11:37.997}
\end{figure}

小さい\(a\)では、\(a^{-4}\)が\(a^{-3}\)より速く増大するため放射項が勝つ：
\begin{equation}
a\ll a_{\rm eq}
\quad\Longrightarrow\quad
H^2\simeq\frac{C_r}{a^4},
\qquad
a(t)\propto t^{1/2}.
\end{equation}
より大きい\(a\)では、物質が勝つ：
\begin{equation}
a\gg a_{\rm eq}
\quad\Longrightarrow\quad
H^2\simeq\frac{C_m}{a^3},
\qquad
a(t)\propto t^{2/3}.
\end{equation}
等しいスケールは、2つの項を等しいと置くことから直接従う：
\begin{equation}
a_{\rm eq}=\frac{C_r}{C_m}.
\label{eq:c02-equality-scale}
\end{equation}
この等しいスケールは光子の脱結合と同じ事象ではない。等しいスケールは2つのエネルギー密度がいつ一致するかを問う；脱結合は相互作用がいつ物質と放射を熱的接触に保つのをやめるかを問う。

放射部門には光子と、超相対論的に運動するあらゆる種が含まれる。非常に軽いニュートリノは十分高温でこのように振る舞い、重力子の背景も同じ仕方でスケールするだろう。いったん質量のある種が非相対論的になると、その支配的な静止質量寄与が物質のようにスケールし始める。冷たい暗黒物質は放射項ではなく\(C_m/a^3\)に属する。

\begin{classroomqa}
\classquestion{方程式~\eqref{eq:c02-mixed-friedmann}のどの項が最初に優勢になり、どれが後に優勢になるか？}
\lecturerresponse{小さい\(a\)では、\(a^{-4}\)が\(a^{-3}\)を追い越すので、放射が初期膨張を支配し\(a\propto t^{1/2}\)を与える。\(a\)が増大するにつれて、放射対物質比は\(1/a\)として減少する；すると物質が膨張を支配し\(a\propto t^{2/3}\)を与える。\lecturetimestamp{01:11:46--01:13:38}}
\end{classroomqa}

\begin{classroomqa}
\audiencequestion{物質と放射の相対的な量は一定であると仮定されているのですか。}
\lecturerresponse{それらの密度は一定ではないが、各成分が効率的にエネルギーを交換しなくなり、粒子の生成や消滅が無視できるようになった後は、係数\(C_m\)と\(C_r\)を一定とみなすことができる。そのとき物質は保存された共動静止質量を、自由な放射は保存された共動光子数を担い、赤方偏移がその余分な因子\(a^{-1}\)を与える。強い相互作用や粒子種の消滅が起こる時代には、この収支計算を精密化する必要がある。\lecturetimestamp{01:13:39--01:14:53}}
\end{classroomqa}

\begin{classroomqa}
\audiencequestion{物質と放射はいつ強い相互作用をやめたのか。}
\lecturerresponse{粗い桁の見積もりでは、光子が通常の物質から脱結合するまでに\(10^5\)年である。正確な再結合の歴史は後の計算であり、物質--放射等量と混同してはならない。ここで必要な点は、脱結合後には二つの成分を近似的に別々に追跡してよいということだけである。\lecturetimestamp{01:27:29--01:27:46}}
\end{classroomqa}

\begin{classroomqa}
\audiencequestion{膨張をたどる銀河が存在する前に、初期の\(a\propto t^{1/2}\)の時代があったことを、どうして知ることができるのですか。}
\lecturerresponse{我々は、膨張の歴史が残す物理的な記録からそれを推論する。原始元素合成は反応率と\(H(t)\)の初期の競合に依存し、宇宙マイクロ波背景放射はより後の熱的状態と摂動を記録している。銀河の観測は物質時代の振る舞いをより直接的にする一方、放射時代はこの初期宇宙の証拠から再構成される。\lecturetimestamp{01:27:46--01:30:10}}
\end{classroomqa}

物質と放射のモデルは意図的に不完全である。それは宇宙が放射優勢から物質優勢へと移行する古典的な枠組みである。現在の加速膨張には別の源、暗黒エネルギーが必要であり、それは質的に異なる項を加えることになる。

\section{赤方偏移したエネルギーの行き先}

\(E_\gamma\propto a^{-1}\)の減少は避けられない問いを提起する。光子のエネルギーはどこへ行ったのか。反射壁をもつ文字通りの膨張する箱では、放射は圧力を及ぼし、壁に対して仕事をする。熱力学的な記述は
\begin{equation}
dE=-p\,dV.
\label{eq:c02-pdv}
\end{equation}
である。\(V\)が増加すると、放射のエネルギーは減少する。失われたエネルギーは、壁を動かすものの運動エネルギー、ひずみ、または熱として現れうる。

共動箱には物質的な壁がない。それは平均的な膨張とともに動く想像上の境界である。光子はそれを横切るが、一様な放射浴では、入ってくる統計的な流束が去っていく流束と釣り合う。この箱は容器ではないにもかかわらず、有効な収支体積であり続ける。圧力仕事の法則は、共動体積内のエネルギーがどのように発展するかを依然として記述している。

放射に対して\(p=\rho/3\)であり、\(E=\rho V\)である。式~\eqref{eq:c02-pdv}はそのとき
\begin{align}
d(\rho V)&=-\frac{\rho}{3}\,dV,\\
V\,d\rho&=-\frac{4}{3}\rho\,dV,\\
\frac{d\rho}{\rho}&=-\frac{4}{3}\frac{dV}{V}.
\end{align}
を与える。\(V\propto a^3\)であるから、積分すると\(\rho\propto a^{-4}\)が得られ、波長の議論と一致する。この熱力学的な導出はまた、光子の数密度だけでなく、各光子のエネルギーが減少することを明らかにする。

\begin{classroomqa}
\audiencequestion{赤方偏移した光子のエネルギーはどこへ行くのか。宇宙の全エネルギーは保存されるのか。}
\lecturerresponse{物理的な壁があれば、放射はそれらに対して仕事をするので、移動は通常の力学的な収支である。宇宙全体については、問いはより微妙である。大域的なエネルギー保存は通常、時間並進対称性から従うが、膨張する時空は一般に時間並進の下で不変ではなく、一意な保存された全エネルギーをもつ必要もない。膨張のエネルギーを含む有用なニュートン的な類推を構築することはできるが、決定的な大域的記述には一般相対性理論が必要である。ここではどちらの記述も有用であるが、ニュートン的モデルから最終的な大域的判定は導かれない。\lecturetimestamp{01:16:51--01:20:58}}
\end{classroomqa}

\begin{classroomqa}
\audiencequestion{共動箱が想像上であり、光子がその境界を横切るのなら、なぜ圧力仕事が依然として有用な描像でありうるのですか。}
\lecturerresponse{一様性が欠けている収支を補う。平均すると、ある共動体積を去るすべての光子は、隣から入ってくる光子によって統計的に置き換えられる。局所的な流束は平均で相殺される一方、物理的な体積と波長は変化する。反射箱は仕事の機構を具体的にする。共動箱は、文字通りの壁なしに、同じ平均化された保存則を担っている。\lecturetimestamp{01:24:22--01:27:28}}
\end{classroomqa}

\begin{classroomqa}
\audiencequestion{膨張は光子密度を減らすだけなのか、それとも個々の光子のエネルギーを減らすのか。}
\lecturerresponse{両方が起こる。単位物理体積あたりの数は\(a^{-3}\)として減少し、宇宙論的赤方偏移は各光子のエネルギーをさらに\(a^{-1}\)の因子だけ下げる。これは断熱冷却の輻射版である。通常の気体も、ゆっくりした膨張の間に仕事をしながら冷却するが、突然の自由膨張は別の熱力学的過程である。\lecturetimestamp{01:42:30--01:45:55}}
\end{classroomqa}

\section{後退、地平線、そして模型の限界}

ハッブル則は、選ばれた宇宙時間のスライス上で固有距離\(D\)にある天体に後退率
\begin{equation}
v_{\rm rec}=HD
\label{eq:c02-recession}
\end{equation}
を割り当てる。十分に大きな\(D\)に対して、式~\eqref{eq:c02-recession}は\(c\)を超えることができる。これは特殊相対性理論に反しない。なぜなら\(v_{\rm rec}\)は、一つの局所慣性系で測定された近傍天体の速度ではないからである。それは、変化する幾何学における広く離れた共動世界線を比較するものである。どの局所観測者も光を速度\(c\)で測定し、いかなる局所的な物質天体も別の天体を光より速く追い越さない。

\begin{classroomqa}
\audiencequestion{十分に遠い銀河は光より速く後退できるか。}
\lecturerresponse{はい、その宇宙論的後退率\(HD\)は\(c\)を超えることができる。そのとき、いかなる局所天体も隣接する観測者を光より速く追い越さず、その結果を超光速シグナル伝達に用いることはできない。それは、遠く離れた共動世界線の間の変化する固有距離についての言明であり、一つの事象で測定された通常の特殊相対論的速度ではない。\lecturetimestamp{01:30:17--01:32:13}}
\end{classroomqa}

超光速後退それ自体は、今放出された光がいつか我々に届くかどうかを決めない。それは将来の膨張史全体についての問いである。減速する物質支配模型では、現在光より速く後退している領域も、関連する地平線が進化するため後に観測可能になりうる。加速膨張が持続すれば、事象の地平線が十分に遠い事象を永久に隠すかもしれない。

\begin{classroomqa}
\audiencequestion{銀河は後退率が光速を超えると消えてしまうのか。}
\lecturerresponse{単純にそうではない。ハッブル半径の横断と真の事象の地平線は異なる概念である。減速模型では、現在超光速で後退している領域からの光が後で我々に届くかもしれない。持続的な暗黒エネルギー支配は大域的な因果構造を変え、究極の地平線を生みうる。どのシグナルが届くかを決めるには時空幾何学が必要であり、\(HD\)のスナップショットでは決まらない。\lecturetimestamp{01:32:14--01:34:18}}
\end{classroomqa}

特殊相対論的な言葉は、大域的に適用すると誤解を招くこともある。粒子の不変静止質量は速度とともに増加しない。そのエネルギーと運動量は観測者の局所系に依存する。質量をもつ粒子は局所的な光速を超えて加速することはできないが、遠くの共動銀河の後退は、それが単一の局所系を通過することではない。

\begin{classroomqa}
\audiencequestion{急速に後退する銀河はより大きな相対論的質量を獲得するのか。}
\lecturerresponse{``相対論的質量''という言葉を使わない方が明快である。銀河は不変静止質量をもち、そのエネルギーと運動量は系に依存する。さらに、宇宙論的距離をまたいで割り当てられた後退率は、特殊相対論の式に現れる局所速度ではない。これらの概念は、遠方の事象を時空計量を用いて比較した後に初めて一致する。\lecturetimestamp{01:38:34--01:40:05}}
\end{classroomqa}

\begin{classroomqa}
\audiencequestion{遠方の後退率が超光速であるとき、時間軸と空間軸はどうなるのか。}
\lecturerresponse{その問いを決着させる傾いた軸をもつ単一の大域ローレンツ系は存在しない。特殊相対性理論は平坦な時空における慣性系を比較するが、宇宙論は広がった曲がった時空にわたって局所系を関係づける計量を必要とする。ニュートン的な推論は小さな領域では信頼できるが、大域的な構成は一般相対性理論を待たねばならない。\lecturetimestamp{01:40:07--01:42:30}}
\end{classroomqa}

ニュートン的なエネルギー保存は、物質について正しいスケール因子の方程式を与え、曲率項さえ予期していた。波動力学と熱力学は輻射のスケーリングを与える。しかし、空間幾何学、大域的能量、地平線は局所的なニュートン的問いではなく、以下に続く計量による記述を必要とする。

\section{物質、輻射、曲率}

物質のみの臨界解は
\begin{equation}
\rho_m\propto a^{-3},
\qquad
a(t)\propto t^{2/3}.
\end{equation}
に従う。単位質量あたりのニュートン的エネルギーを復元すると
\begin{equation}
H^2=\frac{8\pi G}{3}\rho_m+\frac{C}{a^2}.
\end{equation}
が得られる。正の\(C\)は晩期に\(a\sim t\)となる永遠の膨張を与え、\(C=0\)は臨界脱出を与え、負の\(C\)は最大スケール因子とそれに続く再収縮を与える。\(k=-C\)を定義すると、同じ項が相対論的フリードマン方程式の幾何学的側に置かれる。

放射は物質とは異なる。膨張が数密度と量子あたりのエネルギーの両方を変えるからである：
\begin{equation}
\begin{aligned}
\lambda&\propto a,
&\qquad E_\gamma&\propto a^{-1},\\
\rho_r&\propto a^{-4},
& a(t)&\propto t^{1/2}.
\end{aligned}
\end{equation}
したがって、二成分方程式
\begin{equation}
H^2=\frac{C_m}{a^3}+\frac{C_r}{a^4}
\end{equation}
は、放射優勢の初期宇宙に続く物質優勢時代を記述する。暗黒エネルギー、空間幾何、および地平線の因果構造はこのモデルの外にあり、まさにニュートン近似が問題を一般相対論に引き渡すよう指示する場所である。
